\section{Introduction}
\label{sec:intro}

Large language models encode a remarkable range of concepts in their internal representations. Sentiment, truthfulness, and refusal intent can all be recovered from hidden states by training simple linear probes~\cite{marks2023geometry, li2023inference, arditi2024refusal}. The \emph{linear representation hypothesis}~\cite{park2023linear} formalizes this observation: most task-relevant features correspond to linear directions in activation space. But what happens when a feature is \emph{not} linear? Recent work shows that refusal behavior---one of the most safety-critical properties of aligned models---is better described by nonlinear, multi-dimensional structure~\cite{hildebrandt2025refusal, engels2025not}. If models can ``see'' their own linear representations but are blind to their nonlinear ones, this has direct consequences for alignment: self-report-based monitoring may have systematic blind spots.

{\bf Why does this matter?} The ability of an LLM to report on its own internal states---what we call \emph{introspective access}---is a prerequisite for many alignment strategies. Techniques ranging from Constitutional AI to debate-based oversight assume that models can faithfully articulate their reasoning and intentions. If introspective access depends on the geometry of the underlying representation, then some model behaviors may be fundamentally opaque to the model itself. This creates a structural analogy to the conscious/unconscious divide in human cognition~\cite{dehaene2011experimental}: linearly encoded information is ``accessible'' to the model's self-report mechanisms, while nonlinearly encoded information is not.

{\bf What is missing?} Despite substantial progress on both fronts---representation geometry and model self-knowledge---no prior work has directly tested whether these two properties are linked. Studies of linear representations~\cite{marks2023geometry, park2023linear, burns2022discovering} and nonlinear representations~\cite{hildebrandt2025refusal, engels2025not} each characterize feature geometry, but do not ask whether models can introspect differentially on features of different geometry. Conversely, work on factual self-awareness~\cite{tamoyan2025factual} and error acknowledgment~\cite{yang2025when} demonstrates that LLMs possess some introspective capabilities, but does not connect these to the linearity of the underlying feature.

{\bf Our approach.} We propose the first empirical test of the \emph{Introspection--Linearity Hypothesis}: that LLMs' introspective accuracy on a feature is predicted by the linearity of that feature's representation. We employ a two-pronged experimental design. First, we extract residual-stream activations from \qwen and train both linear (logistic regression) and nonlinear (MLP) probes to quantify the geometry of three feature types: sentiment polarity, binary refusal, and harmful-content sub-categories. Second, we administer a battery of introspection tests to \gptfour, measuring self-report accuracy, confidence calibration, and meta-cognitive awareness on tasks spanning the linearity spectrum. We find a consistent introspection gap: 100\% self-report accuracy for clearly linear features versus 90.0--97.5\% for tasks involving nonlinear structure, with significantly miscalibrated confidence at decision boundaries ($p = 0.021$).

Our contributions are as follows:
\begin{itemize}[leftmargin=*,itemsep=0pt,topsep=0pt]
    \item We formalize and test the \emph{Introspection--Linearity Hypothesis}, the first empirical link between representation geometry and introspective access in LLMs.
    \item We demonstrate that features with linear representations (sentiment, binary refusal) are perfectly introspectable (100\% self-report accuracy), while features with nonlinear structure show degraded accuracy (90.0--97.5\%).
    \item We show that confidence calibration tracks representation geometry: models report significantly lower confidence on tasks near nonlinear decision boundaries, revealing partial but imperfect meta-cognitive awareness.
    \item We discuss implications for alignment monitoring, identifying decision boundaries as the primary locus of introspective failure.
\end{itemize}

The rest of this paper is organized as follows. \Secref{sec:related} reviews related work on representation geometry and model self-knowledge. \Secref{sec:method} describes our experimental methodology. \Secref{sec:results} presents results. \Secref{sec:discussion} discusses implications and limitations. \Secref{sec:conclusion} concludes.
